\section{Points of degree 5 on intermediate curves: proof of Theorem B}\label{sec:quintic-intermediate}

\subsection{Reduction to finitely many curves}
\begin{lemma}[Projection]\label{lem:projection}
Let $C\to D$ be a finite morphism over $\Q$. If $D$ has only finitely many points of degree $1$ and of degree
$5$, then $C$ has only finitely many points of degree $5$.
\end{lemma}
\begin{proof}
The image of a point of degree $5$ has degree dividing $5$, and fibres are finite.
\end{proof}
Forgetful maps $X_\Delta(N)\to X_{\Delta'}(M)$, for $M\mid N$ and $\Delta$ mapping into $\Delta'$, are defined
over $\Q$; in particular every $X_\Delta(N)$ maps to $X_0(N)$. By Theorem~A and \cite[Theorem~1.10]{DHJO},
$X_0(N)$ has finitely many points of degree $5$ unless
\[
 N\in\{1,\dots,45\}\cup\{47,\dots,58\}\cup\{61,63,64,65,67,68,72,73,75,80,81,91,109,121,125\},
\]
and every $X_0(N)$ with $N$ outside this list has genus $>1$ (otherwise its rational cusp would give a
function of degree $5$, \S\ref{sec:method-I}), hence finitely many rational points by Faltings
\cite{Faltings1991}; so these $X_0(N)$ can serve as targets in \cref{lem:projection}. So every strict pair with $N\ge192$, and $807$ of the $1011$
strict pairs with $N\le191$, have finitely many points of degree $5$. The strict pairs with $N\le191$ are
enumerated twice, once from primitive rows $(c,d)$ modulo $N$ up to $\Delta$ and once with Sage's
\texttt{GammaH}; full sets of elements, not orders or genera, are used to identify them. Seven more pairs map to
intermediate curves already shown to be negative (\cref{tab:projection}), and fifteen more, of genus $57$ to
$256$, satisfy $119\mu>120000$, so they have $\gon_{\mathbb C}>10$ and finitely many points of degree $5$ by
\cref{lem:abramovich,lem:frey} and \cite[Proposition~2.2]{DerickxSutherland2017} (code S in \cref{tab:negative}). This leaves $182$ strict pairs.

\subsection{The \texorpdfstring{$43$}{43} positive curves}
Each curve in \cref{tab:positive} has a function of degree $5$ over $\Q$, hence infinitely many points of degree
$5$ (\cref{lem:hilbert}):
\begin{itemize}
\item[L] ($13$ curves, $g\le2$): for a rational cusp $P$, $\ell(5P)=6-g>\ell(4P)$;
\item[K] ($14$ curves, $g=3,4$): in genus $3$, $5$ is not a Weierstrass gap at the cusp $\infty$ (read off
 from the orders of vanishing of a basis of cusp forms); in genus $4$ the curve is not hyperelliptic (the
 multiplication map $\mathrm{Sym}^2H^0(K)\to H^0(2K)$ has rank $3g-3=9$, checked on $q$-expansions up to the
 Sturm bound). Then for a rational cusp $P$ the net $|K-P|$ of degree $5$ is base-point free (a base point $Q$
 would give $h^0(P+Q)=2$), so it maps $C$ birationally onto a plane quintic, and projecting from a point of
 $\PP^2(\Q)$ not on this quintic gives a morphism of degree $5$ to $\PP^1$;
\item[W] ($12$ curves): an explicit divisor $D$ of degree $5$ with $\ell(D)=2$ and no base point, on a canonical
 model certified through the Sturm bound and a Hilbert polynomial comparison;
\item[P] ($4$ curves): the $\Q$-pentagonal curves of \cite[Theorems~1.4--1.5]{Orlic2025}.
\end{itemize}

\subsection{The negative curves}
For the remaining $139$ curves we show finiteness. The arguments are listed by code in \cref{tab:negative},
which also lists the $15$ curves of large index (code S) treated above.

\begin{lemma}[Rank zero]\label{lem:rankzero}
If $C(\Q)\ne\varnothing$, $J(C)(\Q)$ is finite and $C$ has no function of degree $5$ over $\Q$, then $C$ has
finitely many points of degree $5$.
\end{lemma}
\begin{proof}
As in \cref{lem:inject}, each point of degree $5$ is the unique effective divisor in its class, and there are
finitely many classes.
\end{proof}

\emph{Rank zero} ($82$ curves). We show $J(\Q)$ finite from nonvanishing $L$-values (exact modular-symbol
$L$-ratios), using Kato \cite[Corollary~14.3]{Kato2004} as applied to modular Jacobians with nontrivial
character in \cite[\S4]{DerickxSutherland2017}. For $25$ of the $74$ curves of class Z we use $L(J_1(N),1)\ne0$:
the forgetful map $X_1(N)\to X_\Delta(N)$ makes $J_\Delta(N)$ an isogeny quotient of $J_1(N)$, so $J_\Delta(N)(\Q)$ is
finite as well. The absence of a function of degree $5$ comes from
\begin{itemize}
\item[Z] ($74$ curves) Orli\'c's classification: genus $\ge5$ and not in Orli\'c's lists of curves of
 $\Q$-gonality at most $5$, so $\gon_\Q\ge6$ \cite{Orlic2025};
\item[B] ($2$ curves, $35\langle-1,8\rangle$ and $55\langle-1,4\rangle$, of genus $7$ and $9$) the curves are
 bielliptic \cite[Theorem~0.2]{JeonKimSchweizer2020}, so a function of degree $5$ would violate
 \cref{lem:CS}: $g\le 2+4=6$;
\item[T] ($3$ curves) Petri lifting: over $\Fp$ every class of $W^1_5$ is $A+P$ for an explicit rational
 base-point free class $A$ of degree $4$, with Petri tangent dimension $1$; \cref{lem:local} identifies $W^1_5$
 near each such class with the Abel curve $A+C$, so a rational pencil of degree $5$ has a base point;
\item[R] ($3$ curves) a complete sieve of rational divisor classes: an explicit subgroup of $J(\Q)$ generated
 by differences of rational points is shown to be all of $J(\Q)$, and every class of $W^1_5(\Fp)$ in its reduction
 is lifted and checked over $\Q$ (each turns out to be a rational class $D_4+P$ with a base point $P$).
 At level $40$ ($\Delta=\langle-1,3\rangle$, genus $7$) the explicit subgroup is generated by the differences of six
 rational points and has order $4032$. An upper bound for $\#J(\Q)$ comes from the Hecke method of
 \cite[\S4.3]{DEH2021}. Rational torsion is fixed by complex conjugation, and for a prime $\ell\nmid40$ the
 Eichler--Shimura relation shows that $1-T_\ell+\ell\langle\ell\rangle=1-T_\ell+T_\ell^2-T_{\ell^2}$ kills rational
 torsion of order prime to $\ell$. Both operators act on $H_1(X,\Z)$, computed from Manin symbols, and the rational
 torsion of order prime to the primes used embeds into the common kernel of these operators on
 $H_1(X,\Q)/H_1(X,\Z)\cong(\Q/\Z)^{2g}$, a finite group computed by a Smith normal form. The relation at $\ell$
 says nothing about $\ell$-primary torsion, so we compute five such kernels: one with all four primes
 $\ell\in\{3,7,11,13\}$, which bounds the primary parts for primes other than these four, and for each of the
 four primes $\ell$ one with the other three primes, which bounds the $\ell$-primary part. All five kernels are
 isomorphic to $\Z/2\times\Z/24\times\Z/168$, of order $8064=2^7\cdot3^2\cdot7$, so $\#J(\Q)$ divides $8064$.
 Reduction is injective on $J(\Q)$ at odd good primes, and the $\ell$-parts of $\#J(\F_p)$ for
 $p\le31$, $p\nmid40\ell$, give $\#J(\Q)\mid2^6\cdot3^2\cdot7^2=28224$. Hence $\#J(\Q)$ divides
 $\gcd(8064,28224)=4032$, and $J(\Q)$ is the explicit subgroup.

 At level $48$ the two curves ($\Delta=\langle-1,5\rangle$ and $\Delta=\langle-1,7,17\rangle$, both of genus $7$)
 have explicit subgroups generated by cuspidal divisors, of orders $2304$ and $12800$, and these are also the upper
 bounds. Let $\pi\colon X_1(48)\to X=X_\Delta(48)$ be the forgetful map, of degree $4$. The pullback
 $\pi^*\colon J(\Q)\to J_1(48)(\Q)$ is induced on homology by the topological transfer
 $H_1(X,\Z)\to H_1(X_1(48),\Z)$, the sum of the four lifts of a cycle, which we compute on Manin symbols (and check
 that it is integral, that $\pi_*\circ\pi^*=4$, and that it commutes with complex conjugation and with $T_7$,
 $T_{11}$). By \cite[Theorem~4.13]{DEH2021}, $J_1(48)(\Q)_{\mathrm{tors}}$ is the rational cuspidal group, which we
 compute in the same Manin-symbol model. So every rational torsion class of $J$ lies in the subgroup of the Hecke
 kernel (computed as at level $40$, with the primes $7,11,13,17,19$) whose transfer lies in the rational
 cuspidal group of $J_1(48)$. This inverse image has order $2304$ for $\Delta=\langle-1,5\rangle$ and $12800$ for
 $\Delta=\langle-1,7,17\rangle$. The argument does not need $\pi^*$ to be injective.
\end{itemize}

\emph{Positive rank} ($57$ curves). None of these curves has a function of degree $5$ over $\Q$: they have genus
at least $8$ and are not in the lists of curves of $\Q$-gonality at most $5$ \cite{JeonKim2007},
\cite[Theorems~1.3--1.5]{Orlic2025}, so their $\Q$-gonality is at least $6$. The five curves of class F and the
curve of class C are treated directly below. For the other $51$ curves (classes E, E0 and D) we use the
classification of curves with infinitely many points of degree at most $4$
\cite[Theorems~1.4--1.6]{DerickxOrlic2026}: none of these $51$ curves has infinitely many points of degree $\le4$, so
if one of them had infinitely many points of degree $5$, its minimum density degree $\min\delta(X/\Q)$ would be $5$.
We use two results of Kadets and Vogt for a nice curve $X$ of genus $g$ over a number field $k$.
\begin{itemize}
\item \cite[Theorem~1.3]{KadetsVogt2025}: \emph{if $\min\delta(X/k)=d$, $m=\lceil d/2\rceil-1$ and $\varepsilon=3d-1-6m$,
 then either there is a nonconstant morphism $\varphi\colon X\to Y$ of degree at least $2$ with
 $d=\min\delta(Y/k)\cdot\deg\varphi$, or $g\le\max\bigl(\frac{d(d-1)}2+1,\,3m(m-1)+m\varepsilon\bigr)$.}
\item \cite[Proposition~6.3 and Table~1]{KadetsVogt2025}: \emph{if $\min\delta(X/k)=5$, then $X$ is a cover of degree
 $5$ of $\PP^1$ or of an elliptic curve of positive rank, or the normalization of a Debarre--Fahlaoui curve, or
 $g\in\{5,6,7,8\}$.} Here a Debarre--Fahlaoui curve \cite[Definition~5.1]{KadetsVogt2025}, \cite{DebarreFahlaoui1993}
 is a geometrically integral curve on $\mathrm{Sym}^2E$, for an elliptic curve $E/k$ of positive rank, in the
 numerical class $(d+j)H-jF$ with $1\le j\le d$, where $H$ is the class of the divisors $\{x+y:x=x_0\}$ and $F$ the
 class of the fibres of the addition map $\mathrm{Sym}^2E\to E$.
\end{itemize}
For $d=5$ we have $m=2$ and $\varepsilon=2$, so the bound in Theorem~1.3 is $g\le\max(11,10)=11$; and since $5$ is
prime, a morphism as in its first alternative has degree $5$ to a curve $Y$ with $\min\delta(Y/\Q)=1$, that is, to
$\PP^1$ or to an elliptic curve of positive rank (\cite[Corollary~2.5]{DerickxOrlic2026} states this for
$g\ge12$). So for $g\ge12$ it suffices to exclude these two kinds of maps. For $9\le g\le11$ Proposition~6.3 adds
only Debarre--Fahlaoui curves. On $\mathrm{Sym}^2E$ we have $H^2=H\cdot F=1$, $F^2=0$ and $K=-2H+F$, so a curve of
class $(5+j)H-jF$ has arithmetic genus $11-j(j-1)/2$ by adjunction, and $F\cdot((5+j)H-jF)=5+j$. Its normalization has
genus at most $11-j(j-1)/2$, so $g\ge9$ forces $j\le2$, and the addition map gives a morphism $X\to E$ of degree $6$
or $7$ to an elliptic curve of positive rank. The curves of genus $8$ ($53\langle-1,4\rangle$ and
$61\langle-1,4\rangle$) are in class F and do not use this argument.
\begin{itemize}
\item[E] ($43$ curves, genus $\ge12$, so only elliptic maps of degree $5$ need to be excluded) every positive-rank elliptic factor $E$ has conductor $N$, multiplicity one
 and odd analytic rank. For such $E$ and $N\le800$, the composition $X_\Delta(N)\to X_0(N)\to E$ with the optimal
 parametrization of $X_0(N)$ is a strong Weil parametrization of $X_\Delta(N)$, of degree $t\,d_E$, where
 $t=[(\Z/N\Z)^\times:\Delta]$ and $d_E$ is the modular degree \cite[Corollary~4.10]{DerickxOrlic2026}. Every morphism
 from $X_\Delta(N)$ to an elliptic curve $E'$ isogenous to $E$ factors, up to translation, through this strong Weil
 parametrization, as in the proof of \cite[Corollary~4.12]{DerickxOrlic2026}, so its degree is divisible by
 $t\,d_E$. In all these cases $t,d_E>1$ (\cref{tab:classE}), so $t\,d_E$ is composite and the degree is not $5$.
 At levels $67$ and $73$ there is no elliptic curve of conductor $N$ and odd analytic rank, and the curves of
 conductor $1$ or $N$ have analytic rank $0$, hence rank $0$; so these seven curves have no positive-rank elliptic
 quotient at all and the condition is vacuous. They could equally be listed under E0; we keep the codes of the
 first version;
\item[E0] ($5$ curves of genus $16$ to $64$, at levels $109$ and $125$) the decomposition of the Jacobian has
 no positive-rank elliptic factor, which is certified by a separate computation of the elliptic factors;
\item[D] ($3$ curves of genus $9$ or $10$: $37\langle-1,10\rangle$, $57\langle-1,2\rangle$, $73\langle-1,6\rangle$)
 here a Debarre--Fahlaoui configuration must also be excluded, that is, elliptic maps of degree $5$, $6$ or $7$:
 at level $73$ there is no positive-rank elliptic factor, and at levels $37$ and $57$ every elliptic map of
 positive rank has degree at least $12$ and $8$ respectively \cite[Corollary~4.10]{DerickxOrlic2026};
\item[C] ($1$ curve, $43\langle-1,2\rangle$) the only possibly positive-rank factor is an elliptic curve $A$, so a
 positive-dimensional translate in $\W$ with dense rational points has the form $x+A$ with $x$ rational; its
 reduction modulo $2$ contains the full coset $\bar x+A(\F_2)$, of order $5$, in $\W(\F_2)$. The set $\W(\F_2)$ has
 $2133$ elements, and $J(\F_2)\cong\Z/38\times\Z/1330$ has a unique subgroup of order $5$, which is therefore
 $A(\F_2)$; no coset of it lies in $\W(\F_2)$. So $\W(\Q)$ is finite;
\item[F] ($5$ curves, levels $53$, $61$, $65$) there is a map $f\colon C\to E$ of degree $4$ whose pullback is
 the positive-rank part $A$, and the surface $Z=\{f^*Q+P\}$ of reducible divisors lies in $\W$. At a good odd
 prime $p$, with $B=A(\Fp)$, we check that every full $B$-coset in $\W(\Fp)$ lies in $Z$, that the intersection
 of the corresponding tangent spaces has dimension at most $2$, and that $A\times C\to\Pic^5C$ has injective
 differential at the relevant points. Since $f_*f^*=[4]$ and $p$ is odd, $A$ has a smooth model over $\Z_p$.
 We lift each point of $B$ to some point of $A(\Z_p)$; for the two curves of genus $11$ at level $65$ we use $p=3$,
 which divides $\#A(\F_3)=6$, so these lifts need not be torsion or form a subgroup, but nothing requires this: a
 rational translate of $A$ in $\W$ and the surface $Z$ are both stable under all of $A$, so both lie in
 $\bigcap_b(\W-a_b)$, and \cref{lem:local} puts every rational translate of $A$ in $\W$ into $Z$.
 A point of degree $5$ in $Z$ would be $f^*Q+P$ with a unique, hence rational, point $P$, which is impossible.
\end{itemize}

\begin{proof}[Proof of Theorem B]
The $43$ positive curves, the $139$ curves above, the $7$ further projections, the $15$ curves with large index
and the $807$ projections to $X_0(N)$ account for all $1011$ strict pairs with $N\le191$, and projection to
$X_0(N)$ handles $N\ge192$.
\end{proof}

\begin{proof}[Proof of Corollary E]
The $39$ positive curves other than Orli\'c's four $\Q$-pentagonal curves have $\Q$-gonality at most $4$: those
of genus $g\le4$ have a rational cusp, so $\gon_\Q\le\max(g,2)\le4$ (\S\ref{sec:method-I}), and the twelve
curves of type W are $\Q$-tetragonal \cite[Theorem~1.3]{Orlic2025}. By \cref{lem:hilbert} they have minimum density degree at most $4$. The four
pentagonal curves have genus $6$, $6$, $8$ and $16$, so finitely many rational points by Faltings
\cite{Faltings1991}, and none of them occurs in the lists of \cite[Theorems~1.4, 1.5 and~1.6]{DerickxOrlic2026}
of curves with infinitely many points of degree $2$, $3$ or $4$. By Theorem~B they have infinitely many points
of degree $5$.
\end{proof}

Every case has a certificate or a precise citation in the accompanying repository
(\texttt{certificates/\allowbreak B\_intermediate\_quintic}), which also contains the full case list.
